
\documentclass{article}
\usepackage{colortbl}
\usepackage{makecell}
\usepackage{multirow}
\usepackage{supertabular}

\begin{document}

\newcounter{utterance}

\centering \large Interaction Transcript for game `eqbench', experiment `eq\_bench', episode 84 with reflexion{-}Qwen3.5{-}27B.
\vspace{24pt}

{ \footnotesize  \setcounter{utterance}{1}
\setlength{\tabcolsep}{0pt}
\begin{supertabular}{c@{$\;$}|p{.15\linewidth}@{}p{.15\linewidth}p{.15\linewidth}p{.15\linewidth}p{.15\linewidth}p{.15\linewidth}}
    \# & \multicolumn{2}{c}{Player} && \multicolumn{2}{c}{Game Master} \\
    \hline

    \theutterance \stepcounter{utterance}  
    & & & \multicolumn{4}{p{0.6\linewidth}}{
        \cellcolor[rgb]{0.9,0.9,0.9}{
            \makecell[{{p{\linewidth}}}]{
                \texttt{\tiny{[P1$\langle$GM]}}
                \texttt{You have attempted this same task before and did not succeed. These are the plans you wrote after those attempts.} \\
\texttt{Trial \#1: New plan: My previous attempt failed because I assigned a high score to "Threatened" (7) and a non{-}zero score to "Resigned" (2), which likely contradicted the specific emotional arc of the dialogue. Charlie's responses ("Don't get smart with me," "That's a low blow," "You don't know a damn thing about me") are defensive and aggressive, indicating he feels personally attacked and his authority is being challenged, but he is actively fighting back rather than feeling helpless (Threatened) or giving up (Resigned). The dialogue ends with Charlie issuing a command ("focus on following protocol"), which is an assertion of control, not resignation. Next time, I will analyze the final lines to determine the dominant emotion: Charlie is clearly **Angry** due to the personal insults and the challenge to his experience. I will assign a high score to Angry (e.g., 9 or 10) and strictly assign 0 to Hopeful, Resigned, and Threatened, as his tone is confrontational and authoritative, not fearful or defeated. I will ensure the scores reflect that his primary reaction is active hostility rather than passive fear or acceptance.} \\
\\ 
\texttt{Trial \#2: New plan: My previous attempt failed because I assigned a score of 0 to "Threatened," overlooking the subtext of Charlie's defensiveness. While Charlie is clearly **Angry** (high score), Mari's specific attacks on his age ("washed up"), his drinking habits, and his competence ("You don't know a damn thing about me") directly challenge his professional identity and authority. This triggers a feeling of being **Threatened** regarding his status and relevance, which is distinct from simple anger. By assigning 0 to Threatened, I missed a key emotional layer. Next time, I will assign a high score to Angry (e.g., 9) and a moderate score to Threatened (e.g., 6 or 7) to reflect the personal nature of the attacks, while keeping Hopeful and Resigned at 0. I will ensure the scores capture both the active hostility and the underlying insecurity provoked by the dialogue.} \\
\\ 
\texttt{Trial \#3: New plan: My previous attempts failed because I inconsistently evaluated the "Threatened" emotion, first assigning it a high score (7) and then justifying a score of 0, leading to a loss in both cases. This suggests the evaluation criteria for "Threatened" is the critical differentiator. Upon re{-}evaluating the dialogue, Charlie's responses are aggressive assertions of authority ("Don't get smart," "focus on following protocol") rather than signs of fear or vulnerability. While Mari attacks his competence, Charlie's reaction is to double down on his power, not to feel threatened. Therefore, the error in the first attempt was likely over{-}scoring "Threatened," and the error in the second was the inconsistency in reasoning. The correct strategy is to recognize that Charlie is purely **Angry** and defensive, with no genuine sense of threat or resignation. I will assign a maximum score to **Angry** (10) and strictly assign **0** to Hopeful, Resigned, and Threatened, ensuring the output reflects a character who is hostile and in control, not fearful.} \\
\\ 
\\ 
\\ 
\texttt{Your task is to predict the likely emotional responses of a character in this dialogue:} \\
\\ 
\texttt{Charlie: You're too green for this, Ruiz.} \\
\texttt{Mari: And you're too washed up, Redburn.} \\
\texttt{Charlie: I've been on this beat longer than you've been alive, kid.} \\
\texttt{Mari: And yet, here we are, stuck on the same case. Seems like experience isn't everything.} \\
\texttt{Charlie: Don't get smart with me, Mari.} \\
\texttt{Mari: Someone has to be the smart one here.} \\
\texttt{Charlie: You think I don't know what you're doing? Trying to prove yourself, make a name?} \\
\texttt{Mari: Better than drinking myself into oblivion.} \\
\texttt{Charlie: That's a low blow, Ruiz.} \\
\texttt{Mari: Truth hurts, doesn't it?} \\
\texttt{Charlie: You don't know a damn thing about me.} \\
\texttt{Mari: And you don't know anything about me, either. So maybe we should just focus on the case.} \\
\texttt{Charlie: Maybe you should focus on following protocol for once.} \\
\texttt{Mari: And maybe you should try embracing the 21st century.} \\
\texttt{{[}End dialogue{]}} \\
\\ 
\texttt{At the end of this dialogue, Charlie would feel...} \\
\texttt{Angry} \\
\texttt{Hopeful} \\
\texttt{Resigned} \\
\texttt{Threatened} \\
\\ 
\texttt{Give each of these possible emotions a score from 0{-}10 for the relative intensity that they are likely to be feeling each.} \\
\\ 
\texttt{You must output in the following format, including headings (of course, you should give your own scores), with no additional commentary:} \\
\\ 
\texttt{Angry: <score>} \\
\texttt{Hopeful: <score>} \\
\texttt{Resigned: <score>} \\
\texttt{Threatened: <score>} \\
\\ 
\\ 
\texttt{{[}End of answer{]}} \\
\\ 
\texttt{Remember: zero is a valid score, meaning they are likely not feeling that emotion. You must score at least one emotion > 0.} \\
\\ 
\texttt{Your answer:} \\
            }
        }
    }
     \\ \\

    \theutterance \stepcounter{utterance}  
    & \multicolumn{4}{p{0.6\linewidth}}{
        \cellcolor[rgb]{0.561,0.737,0.561}{
            \makecell[{{p{\linewidth}}}]{
                \texttt{\tiny{[P1$\rangle$GM]}}
                \texttt{Angry: 10} \\
\texttt{Hopeful: 0} \\
\texttt{Resigned: 0} \\
\texttt{Threatened: 0} \\
            }
        }
    }
    & & \\ \\

    \theutterance \stepcounter{utterance}  
    & & & \multicolumn{2}{p{0.3\linewidth}}{
        \cellcolor[rgb]{0.9,0.9,0.9}{
            \makecell[{{p{\linewidth}}}]{
                \texttt{\tiny{[GM$|$GM]}}
                \texttt{Angry: 10} \\
\texttt{Hopeful: 0} \\
\texttt{Resigned: 0} \\
\texttt{Threatened: 0} \\
            }
        }
    }
    & & \\ \\

    \theutterance \stepcounter{utterance}  
    & & & \multicolumn{2}{p{0.3\linewidth}}{
        \cellcolor[rgb]{0.9,0.9,0.9}{
            \makecell[{{p{\linewidth}}}]{
                \texttt{\tiny{[GM$|$GM]}}
                \texttt{Angry: 8} \\
\texttt{Hopeful: 0} \\
\texttt{Resigned: 4} \\
\texttt{Threatened: 6} \\
            }
        }
    }
    & & \\ \\

    \theutterance \stepcounter{utterance}  
    & & & \multicolumn{2}{p{0.3\linewidth}}{
        \cellcolor[rgb]{0.9,0.9,0.9}{
            \makecell[{{p{\linewidth}}}]{
                \texttt{\tiny{[GM$|$GM]}}
                \texttt{game\_result = LOSE} \\
            }
        }
    }
    & & \\ \\

    \theutterance \stepcounter{utterance}  
    & & \multicolumn{4}{p{0.6\linewidth}}{
        \cellcolor[rgb]{0.95,0.95,0.95}{
            \makecell[{{p{\linewidth}}}]{
                \texttt{\tiny{[P1$\leftrightarrow$P]}}
                \texttt{Post{-}episode Reflexion (not sent to GM):} \\
\texttt{New plan: My previous attempts failed because I oscillated between over{-}scoring "Threatened" and under{-}scoring it, failing to recognize that the dialogue's conclusion specifically targets Charlie's sense of professional obsolescence. While Charlie is undeniably **Angry** (score 10), Mari's final line ("embrace the 21st century") and the preceding attacks on his age and drinking habits directly challenge his relevance and authority. This specific type of attack—questioning one's competence and place in the world—induces a feeling of being **Threatened** (score 6{-}7), distinct from simple anger. The error in my last attempt was assigning 0 to Threatened, ignoring the subtext of his defensive posturing. Next time, I will assign a high score to **Angry** (10) and a significant, non{-}zero score to **Threatened** (e.g., 7) to reflect the personal and professional insecurity Mari has triggered, while keeping **Hopeful** and **Resigned** at 0, as Charlie is actively fighting back rather than giving up or looking forward.} \\
            }
        }
    }
    & \\ \\

\end{supertabular}
}

\end{document}
